%% macros.tex

\newcommand{\R}{{\mathbb{R}}}
\newcommand{\N}{{\mathbb{N}}}

\newcommand{\set}[1]{{\left\{#1\right\}}}
\newcommand{\given}{\;\middle|\;}

\NewDocumentCommand{\pbrac}{om}{{\IfNoValueTF{#1}{\left\lparen#2\right\rparen}{#1\lparen#2#1\rparen}}}
\NewDocumentCommand{\bbrac}{om}{{\IfNoValueTF{#1}{\left\lbrack#2\right\rbrack}{#1\lbrack#2#1\rbrack}}}

\NewDocumentCommand{\abs}{om}{{\IfNoValueTF{#1}{\left\lvert#2\right\rvert}{#1\lvert#2#1\rvert}}}
\NewDocumentCommand{\norm}{om}{{\IfNoValueTF{#1}{\left\lVert#2\right\rVert}{#1\lVert#2#1\rVert}}}

\newcommand{\dotp}[2]{{{#1}^\top#2}}
\newcommand{\outp}[2]{{#1{#2}^\top}}
\ExplSyntaxOn{}
\NewDocumentCommand{\qf}{mmo}{
    \IfValueTF{#3}{
        \str_if_eq:nnT {#1} {#3} {
            \PackageWarning{macros}{Argument 3 not needed}
        }
        {#1}^\top#2\,#3
    }{
        {#1}^\top#2\,#1
    }
}
\ExplSyntaxOff{}

\newcommand{\inv}[1]{{{#1}^{-1}}}

\newcommand{\func}[2]{{#1\!\left\lparen#2\right\rparen}}

\LoopCommands\lettersUppercase[cal#1]{\newmathcommand#2{{\mathcal{#1}}}}
\newcommand{\classC}{{\calC}}

\LoopCommands{{argmin}{argmax}}[#1]{\newmathcommand#2{\mathop{\mathrm{#1}}}}
\LoopCommands{{minimize}{maximize}}[#1]{\DeclareMathOperator*{#2}{#1}}
\DeclareMathOperator{\subjectto}{subject~to}
\newcommand{\smalloh}[1]{\func{{\scriptstyle\calO}}{#1}}
\newcommand{\bigoh}[1]{\func{\calO}{#1}}

\LoopCommands{xyzqfk}[#1]{\declaremathcommandPIE#2{{\mathbf{#1}}##1##2\IfEmptyTF{##3}{}{^{(\GetExponent{##3})}}}}
\LoopCommands\lettersGreekAll[bf#1]{\newmathcommand#2{{\boldsymbol{#3}}}}

\LoopCommands\lettersAll[bf#1]{\newmathcommand#2{{\mathbf{#1}}}}
\LoopCommands{IAMKBDF}[#1]{\newmathcommand#2{{\mathbf{#1}}}}

\newcommand{\half}{{\frac{1}{2}}}
\newcommand{\zero}{{\mathbf{0}}}
